% =========================================================
% Defining Rational Arithmetic from Representatives
% =========================================================

\section{Defining Rational Arithmetic from Representatives}
\label{sec:rational-arithmetic-from-representatives}

This section begins where the quotient construction leaves off: operations are first written using representatives, and then their well-definedness is proved so that they descend to genuine operations on rational numbers.

\begin{definition}[Addition and multiplication]
\label{def:rational-operations}

Let $[(a,b)],[(c,d)]\in\mathbb{Q}$, where $(a,b),(c,d)\in W$.  Define
\[
[(a,b)] + [(c,d)] := [(ad+bc,bd)]
\]
and
\[
[(a,b)]\cdot[(c,d)] := [(ac,bd)].
\]

\end{definition}

\begin{lemma}[Addition Is Well-Defined on Rational Classes]
\label{lem:rational-addition-well-defined}
Let $a,b,c,d,e,f,g,h\in\mathbb{Z}$ with $bdfh\neq 0$. If
\[
(a,b)\sim(e,f)
\qquad\text{and}\qquad
(c,d)\sim(g,h),
\]
then
\[
(ad+bc,bd)\sim(eh+fg,fh).
\]
\end{lemma}

\begin{remark*}[Proof]
\hyperref[prf:rational-addition-well-defined]{Go to proof.}
\end{remark*}

\begin{lemma}[Multiplication Is Well-Defined on Rational Classes]
\label{lem:rational-multiplication-well-defined}
Let $a,b,c,d,e,f,g,h\in\mathbb{Z}$ with $bdfh\neq 0$. If
\[
(a,b)\sim(e,f)
\qquad\text{and}\qquad
(c,d)\sim(g,h),
\]
then
\[
(ac,bd)\sim(eg,fh).
\]
\end{lemma}

\begin{remark*}[Proof]
\hyperref[prf:rational-multiplication-well-defined]{Go to proof.}
\end{remark*}

\begin{remark}[Why well-definedness matters]
\label{rem:well-definedness-matters}

These formulas are written in terms of representatives.  Because one rational number has many
representatives, it must be proved that choosing different representatives leads to the same
resulting equivalence class.  This is the first serious instance of a very common pattern in
mathematics:
\[
\text{representatives}
\longrightarrow
\text{operations}
\longrightarrow
\text{well-defined operations on equivalence classes}.
\]

\end{remark}

\begin{remark}[Three levels of representation]
\label{rem:rational-three-levels}

When working with rational numbers, it is essential to distinguish among the following three objects.

\begin{center}
\begin{tabular}{lll}
\toprule
\textbf{Object} & \textbf{Lives in} & \textbf{Meaning} \\
\midrule
$(a,b)$ & $W$ & raw integer pair \\
$[(a,b)]$ & $\mathbb{Q}$ & rational number \\
$\dfrac{a}{b}$ & notation & informal representation of $[(a,b)]$ \\
\bottomrule
\end{tabular}
\end{center}

The pair $(a,b)$ is not itself a rational number.  The rational number is the entire equivalence
class
\[
[(a,b)] = \{(c,d)\in W : ad = bc\}.
\]

\end{remark}

\begin{remark}[A rational number is an infinite set]
\label{rem:rational-infinite-set}

For every nonzero integer $k$,
\[
(ka,kb)\sim(a,b).
\]
Hence the class $[(a,b)]$ contains infinitely many representatives:
\[
(a,b),(2a,2b),(3a,3b),\dots,(-a,-b),(-2a,-2b),\dots
\]
Thus every rational number is an infinite set.
\end{remark}

\begin{remark}[Canonical representative]
\label{rem:reduced-rational}

Although a rational number has infinitely many representatives, each nonzero rational number can be
written uniquely in the form
\[
\frac{a}{b}
\qquad\text{with}\qquad
\gcd(a,b)=1,
\quad b>0.
\]
This is the \emph{reduced form} of the rational number.

\end{remark}

\subsection{Consequences}

\begin{remark}[Quotient structures]
\label{rem:quotient-pattern}

The construction of $\mathbb{Q}$ is the first major example in these notes of the general pattern
\[
\text{set} + \text{equivalence relation} \longrightarrow \text{quotient structure}.
\]
The same pattern later appears in modular arithmetic, quotient groups, quotient rings, quotient
spaces, and projective geometry.

\end{remark}

\subsection{Algebra of Equivalence Classes}

\begin{remark}[Role of the equivalence-class arithmetic theorems]
\label{rem:ec-arithmetic-role}

The theorems in this subsection work entirely at the level of
equivalence classes and integer arithmetic.  They require nothing
beyond \hyperref[def:rational-operations]{Definition~\ref*{def:rational-operations}},
the well-definedness lemmas, and facts about $\mathbb{Z}$.

In particular, \hyperref[thm:ec-negation]{Theorem~\ref*{thm:ec-negation}}
and \hyperref[thm:ec-reciprocal]{Theorem~\ref*{thm:ec-reciprocal}} do
not appeal to field axioms.  They \emph{establish}, by direct
computation, that $[(-a,b)]$ and $[(b,a)]$ satisfy the defining
conditions for additive and multiplicative inverses respectively.
These facts are therefore \emph{inputs to} the proof that $\mathbb{Q}$
is a field, not consequences of it.  The field proof can cite them
rather than repeat the computation.

The later \hyperref[sec:rational-algebra]{rational algebra section}
re-expresses the same identities using fraction notation $a/b$ inside
the already-established field.  Those re-derivations are about
\emph{notation and usability}, not about new mathematical content.

\end{remark}

\begin{theorem}[Integer embedding preserves addition]
\label{thm:ec-integer-embedding-add}
For all $m,n\in\mathbb{Z}$,
\[
[(m,1)] + [(n,1)] = [(m+n,\,1)].
\]
\end{theorem}

\begin{remark*}[Proof]
\hyperref[prf:ec-integer-embedding-add]{Go to proof.}
\end{remark*}

\begin{theorem}[Integer embedding preserves multiplication]
\label{thm:ec-integer-embedding-mul}
For all $m,n\in\mathbb{Z}$,
\[
[(m,1)] \cdot [(n,1)] = [(mn,\,1)].
\]
\end{theorem}

\begin{remark*}[Proof]
\hyperref[prf:ec-integer-embedding-mul]{Go to proof.}
\end{remark*}

\begin{remark}[The integer embedding]
\label{rem:integer-embedding}
These two theorems establish that the map $n\mapsto[(n,1)]$ is a ring
homomorphism from $\mathbb{Z}$ into $\mathbb{Q}$: it sends the integer
operations to the rational operations.  Moreover $[(m,1)]=[(n,1)]$ if
and only if $m=n$, so the map is injective and $\mathbb{Z}$ sits inside
$\mathbb{Q}$ as a subring.  We therefore identify each integer $n$ with
the rational class $[(n,1)]$ and write $n$ for $[(n,1)]$ without further
comment.
\end{remark}

\begin{theorem}[Zero class criterion]
\label{thm:ec-zero-numerator}
For every $[(a,b)]\in\mathbb{Q}$,
\[
[(a,b)] = [(0,1)]
\quad\Longleftrightarrow\quad
a = 0.
\]
\end{theorem}

\begin{remark*}[Proof]
\hyperref[prf:ec-zero-numerator]{Go to proof.}
\end{remark*}

\begin{theorem}[Scaling of equivalence classes]
\label{thm:ec-scaling}
For every $[(a,b)]\in\mathbb{Q}$ and every nonzero integer $e\neq 0$,
\[
[(a,b)] = [(ae,be)].
\]
\end{theorem}

\begin{remark*}[Proof]
\hyperref[prf:ec-scaling]{Go to proof.}
\end{remark*}

\begin{theorem}[Negation of an equivalence class]
\label{thm:ec-negation}
For every $[(a,b)]\in\mathbb{Q}$,
\[
-[(a,b)] = [(-a,b)] = [(a,-b)].
\]
\end{theorem}

\begin{remark*}[Proof]
\hyperref[prf:ec-negation]{Go to proof.}
\end{remark*}

\begin{theorem}[Subtraction of equivalence classes]
\label{thm:ec-subtraction}
For every $[(a,b)],[(c,d)]\in\mathbb{Q}$,
\[
[(a,b)] - [(c,d)] = [(ad-bc,\,bd)].
\]
\end{theorem}

\begin{remark*}[Proof]
\hyperref[prf:ec-subtraction]{Go to proof.}
\end{remark*}

\begin{theorem}[Reciprocal of a nonzero equivalence class]
\label{thm:ec-reciprocal}
For every $[(a,b)]\in\mathbb{Q}$ with $a\neq 0$,
\[
[(a,b)]^{-1} = [(b,a)].
\]
\end{theorem}

\begin{remark*}[Proof]
\hyperref[prf:ec-reciprocal]{Go to proof.}
\end{remark*}
